% =====================================================================
%  Master rad — glavni fajl
%  Build: lualatex → biber → lualatex → lualatex
%  (VS Code: recept "lualatex -> biber -> lualatex x2")
%
%  Izgled i struktura prate odobreni master rad sa ETF-a (master-primjer.pdf):
%  margine 2,54 cm, Times New Roman 12, prored 1,5 (kao u Word-u),
%  uvlačenje prvog reda 1,27 cm, broj stranice dolje u sredini,
%  UVOD i ZAKLJUČAK bez broja, poglavlja 1..n između njih.
% =====================================================================
\documentclass[12pt,a4paper,oneside]{report}

% --- Stranica i font ----------------------------------------------------
\usepackage[margin=2.54cm,footskip=1.25cm]{geometry}
\usepackage{fontspec}
\setmainfont{Times New Roman}
\setmonofont{Consolas}[Scale=0.9]
\usepackage[english,main=serbian]{babel}

% Word "prored 1,5" za TNR 12 = razmak redova 20,7 pt (LaTeX osnovno: 14,5 pt)
\usepackage{setspace}
\setstretch{1.43}
\setlength{\parindent}{1.27cm}
\setlength{\parskip}{6pt plus 2pt}

% Broj stranice dolje u sredini
\usepackage{fancyhdr}
\pagestyle{fancy}
\fancyhf{}
\fancyfoot[C]{\thepage}
\renewcommand{\headrulewidth}{0pt}
\fancypagestyle{plain}{\fancyhf{}\fancyfoot[C]{\thepage}\renewcommand{\headrulewidth}{0pt}}

% --- Naslovi: "1. NASLOV" (14 pt), "1.1. Naslov" (12 pt), svi podebljani --
\usepackage{titlesec}
\titleformat{\chapter}[hang]{\normalfont\fontsize{14}{17}\selectfont\bfseries\raggedright}{\thechapter.}{0.4em}{}
\titleformat{name=\chapter,numberless}[hang]{\normalfont\fontsize{14}{17}\selectfont\bfseries\raggedright}{}{0pt}{}
\titlespacing*{\chapter}{0pt}{0pt}{24pt}
\titleformat{\section}{\normalfont\fontsize{12}{15}\selectfont\bfseries}{\thesection.}{0.4em}{}
\titlespacing*{\section}{0pt}{12pt}{6pt}
\titleformat{\subsection}{\normalfont\fontsize{12}{15}\selectfont\bfseries}{\thesubsection.}{0.4em}{}
\titlespacing*{\subsection}{0pt}{12pt}{6pt}

% --- Sadržaj: tačkice za sve nivoe, poglavlja nisu podebljana ------------
\usepackage{titletoc}
\titlecontents{chapter}[0pt]{\addvspace{2pt}}
  {\thecontentslabel.\ }{}{\titlerule*[0.25em]{.}\contentspage}
\titlecontents{section}[0.6cm]{}
  {\thecontentslabel.\ }{}{\titlerule*[0.25em]{.}\contentspage}
\titlecontents{subsection}[1.2cm]{}
  {\thecontentslabel.\ }{}{\titlerule*[0.25em]{.}\contentspage}
\setcounter{tocdepth}{2}

% --- Slike, tabele, kod --------------------------------------------------
\usepackage{graphicx}
\graphicspath{{slike/}}
\usepackage{xcolor}
\usepackage{booktabs}
\usepackage{tabularx}
\usepackage{array}
\usepackage{longtable}
\usepackage{float}
\usepackage{enumitem}
\usepackage{listings}
% =====================================================================
%  Podešavanje listinga koda (Python i RISC-V asembler)
% =====================================================================
\lstset{
  basicstyle=\ttfamily\footnotesize\singlespacing,
  numbers=left,
  numberstyle=\tiny\color{gray},
  frame=single,
  rulecolor=\color{gray!60},
  breaklines=true,
  columns=fullflexible,
  keepspaces=true,
  showstringspaces=false,
  tabsize=4,
  extendedchars=true,
  keywordstyle=\bfseries\color{blue!60!black},
  commentstyle=\itshape\color{green!40!black},
  stringstyle=\color{red!50!black},
  literate={—}{{---}}1 {→}{{$\rightarrow$}}1,
}

% Podskup RV32I + pseudo-instrukcije + direktive
\lstdefinelanguage{rvasm}{
  sensitive=false,
  morekeywords={add,sub,and,or,xor,sll,srl,sra,slt,sltu,
                addi,andi,ori,xori,slti,sltiu,slli,srli,srai,
                lw,sw,lb,sb,jalr,beq,bne,blt,bge,bltu,bgeu,lui,auipc,jal,
                nop,li,la,mv,j,ret},
  morekeywords=[2]{.text,.data,.word,.byte,.ascii,.asciiz,.space},
  keywordstyle=[2]\color{purple!70!black},
  morecomment=[l]{\#},
  morestring=[b]",
  alsoletter={.},
}

\lstdefinestyle{python}{language=Python}
\lstdefinestyle{asm}{language=rvasm}


% Natpisi: "Slika 4. Opis (izvor)" — kurziv, centrirano, jedinstvena numeracija
\usepackage[labelsep=period,font=it,justification=centering,skip=6pt]{caption}
\counterwithout{figure}{chapter}
\counterwithout{table}{chapter}
\AtBeginDocument{\counterwithout{lstlisting}{chapter}}

% --- Literatura: APA-oblik (Autor, godina), numerisan spisak --------------
\usepackage{csquotes}
\usepackage[backend=biber,style=authoryear,maxcitenames=2,mincitenames=1,
            maxbibnames=20,giveninits=true,uniquelist=false,uniquename=false,
            dashed=false,urldate=long]{biblatex}
\addbibresource{literatura.bib}
% =====================================================================
%  Literatura u APA obliku, kao u master-primjer.pdf:
%   u tekstu:  (Aho et al., 2006), (Patterson & Hennessy, 2020)
%   spisak:    1. Aho, A. V., Lam, M. S., Sethi, R., & Ullman, J. D. (2006). Naslov. Izdavač.
%   podnaslovi spiska ("knjige i članci", "Internet izvori") centrirani, numeracija od 1.
% =====================================================================

% Veznici i skraćenice kao u APA stilu (i u srpskom tekstu primjera)
\DefineBibliographyStrings{serbian}{
  andothers = {et al\adddot},
  and       = {\&},
  urlseen   = {pristupljeno},
  edition   = {izd\adddot},
}
% U tekstu: (Patterson & Hennessy, 2020); u spisku (APA): Patterson, D. A., & Hennessy, J. L.
\DeclareDelimFormat{finalnamedelim}{\addspace\&\space}
\DeclareDelimFormat[bib,biblist]{finalnamedelim}{\addcomma\space\&\space}
\renewcommand*{\nameyeardelim}{\addcomma\space}
\DeclareNameAlias{sortname}{family-given}
\DeclareNameAlias{default}{family-given}

% Godina bez tačke: (2006), ne (2006.) — serbian.lbx dodaje tačku iza godine
\DefineBibliographyExtras{serbian}{%
  \protected\def\mkbibdatelong#1#2#3{%
    \iffieldundef{#3}{}{\mkbibordinal{\thefield{#3}}\iffieldundef{#2}{}{\nobreakspace}}%
    \iffieldundef{#2}{}{\mkbibmonth{\thefield{#2}}\iffieldundef{#1}{}{\space}}%
    \iffieldbibstring{#1}{\bibstring{\thefield{#1}}}{\dateeraprintpre{#1}\mkyearzeros{\thefield{#1}}}}%
  \protected\def\mkbibdateshort#1#2#3{%
    \iffieldundef{#3}{}{\mkbibordinal{\thefield{#3}}\iffieldundef{#2}{}{\bibdatesep}}%
    \iffieldundef{#2}{}{\mkbibordinal{\thefield{#2}}\iffieldundef{#1}{}{\bibdatesep}}%
    \iffieldbibstring{#1}{\bibstring{\thefield{#1}}}{\dateeraprintpre{#1}\mkyearzeros{\thefield{#1}}}}%
}

% Internet izvori: "…, https://…, pristupljeno: 27. oktobar 2026."
\DeclareFieldFormat{url}{\url{#1}}
\AtBeginDocument{\urlstyle{same}}
\DeclareFieldFormat{urldate}{\bibstring{urlseen}\addcolon\space#1}
\renewbibmacro*{url+urldate}{%
  \printfield{url}%
  \iffieldundef{urlyear}{}{\setunit{\addcomma\space}\printurldate}}

% Naslovi knjiga kurzivom, naslovi radova bez navodnika
\DeclareFieldFormat[article,inproceedings,incollection]{title}{#1}

% Numerisani spisak, podnaslovi centrirani
\defbibenvironment{bibliography}
  {\begin{enumerate}[label=\arabic*., leftmargin=0.9cm, labelwidth=0.7cm, labelsep=0.2cm,
                     align=left, itemsep=2pt]}
  {\end{enumerate}}
  {\item}
\defbibheading{subbibliography}[]{\par\begin{center}#1\end{center}}


\usepackage[hidelinks]{hyperref}

% --- Nazivi na srpskom ---------------------------------------------------
\addto\captionsserbian{%
  \renewcommand{\contentsname}{SADRŽAJ}%
  \renewcommand{\bibname}{LITERATURA}%
  \renewcommand{\listfigurename}{SPISAK SLIKA}%
  \renewcommand{\listtablename}{SPISAK TABELA}%
  \renewcommand{\figurename}{Slika}%
  \renewcommand{\tablename}{Tabela}%
}
\renewcommand{\lstlistingname}{Listing}

% Nenumerisano poglavlje koje ipak ulazi u sadržaj (UVOD, ZAKLJUČAK, ...)
\newcommand{\nenumerisano}[1]{\chapter*{#1}\addcontentsline{toc}{chapter}{#1}}

% Prilozi: "Prilog 1. Naslov" (u sadržaju kao sekcija; \ref daje broj priloga)
\newcounter{prilog}
\newcommand{\prilog}[1]{\refstepcounter{prilog}%
  \section*{Prilog \theprilog. #1}\addcontentsline{toc}{section}{Prilog \theprilog. #1}}

% --- Radne napomene (uputstva šta ide u sekciju) -------------------------
% Za finalnu verziju: \napomenefalse — napomene nestaju iz PDF-a.
\newif\ifnapomene
\napomenetrue
\newcommand{\napomena}[1]{\ifnapomene{\par\raggedright\color{gray}\small\itshape [#1]\par}\fi}

% =====================================================================
%  Podaci o radu — mijenjati SAMO ovdje; koriste se na koricama,
%  naslovnoj strani i strani sa informacijama.
% =====================================================================
\newcommand{\RadUniverzitet}{Univerzitet Crne Gore}
\newcommand{\RadFakultet}{[Naziv fakulteta]}
\newcommand{\RadStudijskiProgram}{[Naziv studijskog programa]}
\newcommand{\RadVrsta}{Master rad}

\newcommand{\RadNaslov}{Dizajn i implementacija jednostavnog asemblerskog prevodioca za prostu arhitekturu}
\newcommand{\RadNaslovEn}{Design and Implementation of a Simple Assembler for a Simple Architecture}

\newcommand{\RadAutor}{Milica Delibašić}
\newcommand{\RadMentor}{[akademsko zvanje, ime i prezime mentora]}
\newcommand{\RadMjesto}{Podgorica}
\newcommand{\RadGodina}{[godina]}


% =====================================================================
\begin{document}

% --- Uvodni dio: bez brojeva stranica ------------------------------------
\pagenumbering{gobble}
% Korice i naslovna strana — identične (Smjernice), izgled kao u primjeru ETF-a.
\newcommand{\RadKorice}{%
  \begin{titlepage}
    \setstretch{1.2}
    \centering
    \begin{minipage}[c]{0.24\textwidth}\centering
      \includegraphics[height=2.4cm]{logo-ucg.png}
    \end{minipage}%
    \begin{minipage}[c]{0.5\textwidth}\centering
      \textsc{\RadUniverzitet}\\[0.5em]
      \textsc{\RadFakultet}
    \end{minipage}%
    \begin{minipage}[c]{0.24\textwidth}\centering
      \ifx\RadLogoFakulteta\empty\else\includegraphics[height=2.2cm]{\RadLogoFakulteta}\fi
    \end{minipage}

    \vspace*{\fill}
    {\RadZvanje\ \RadAutor\par}
    \vspace{2.5em}
    {\bfseries\MakeUppercase{\RadNaslov}\par}
    \vspace{1em}
    {MASTER RAD\par}
    \vspace*{\fill}
    \vspace*{\fill}
    {\RadMjesto, \RadGodina. godine\par}
  \end{titlepage}}
\RadKorice
\RadKorice

% Informacije o magistrandu, master radu i komisiji (Smjernice, tačka 2).
\chapter*{Informacije o magistrandu, master radu i komisiji}
\thispagestyle{plain}

\section*{Podaci o magistrandu}
\noindent\begin{tabularx}{\textwidth}{@{}lX@{}}
  Ime i prezime: & \RadAutor \\
  Datum i mjesto rođenja: & [\dots] \\
  Završeni studijski program i godina završetka: & [\dots] \\
\end{tabularx}

\section*{Podaci o master radu}
\noindent\begin{tabularx}{\textwidth}{@{}lX@{}}
  Naziv master studija: & \RadStudijskiProgram \\
  Naslov rada: & \RadNaslov \\
  Fakultet na kojem se rad brani: & \RadFakultet \\
\end{tabularx}

\section*{Podaci o komisiji, prijavi i odbrani}
\noindent\begin{tabularx}{\textwidth}{@{}lX@{}}
  Datum prijave master rada: & [\dots] \\
  Datum sjednice Vijeća na kojoj je prihvaćena tema: & [\dots] \\
  Mentor: & \RadMentor \\
  Komisija za ocjenu rada: & [\dots] \\
  Komisija za odbranu rada: & [\dots] \\
  Lektor: & [\dots] \\
  Datum odbrane: & [\dots] \\
\end{tabularx}

% Izjava o autorstvu — tekst prema obrascu iz primjera.
% Fakultet obično traži potpisan obrazac: tada ovu stranu zamijeniti skeniranim PDF-om
% (\usepackage{pdfpages} i \includepdf{izjava-autorstvo-potpisana.pdf}).
\clearpage
\thispagestyle{empty}
\begingroup
\setlength{\parindent}{0pt}
\vspace*{2cm}
\begin{center}\bfseries Izjava o autorstvu\end{center}
\vspace{2em}

Potpisani/a: \RadAutor\\
Broj indeksa/upisa: \RadIndeks

\begin{center}\bfseries Izjavljujem\end{center}

da je master rad pod nazivom

\begin{center}\bfseries\enquote{\RadNaslov}\end{center}

\begin{itemize}[leftmargin=1.5cm]
  \item rezultat sopstvenog istraživačkog rada,
  \item da predloženi master rad ni u cjelini ni u djelovima nije bio predložen za dobijanje
        bilo koje diplome prema studijskim programima drugih ustanova visokog obrazovanja,
  \item da su rezultati korektno navedeni, i
  \item da nijesam povrijedio/la autorska i druga prava intelektualne svojine koja pripadaju
        trećim licima.
\end{itemize}

\vspace{3em}
U Podgorici, [dd.mm.gggg.] godine

\vspace{2em}
\hfill\begin{tabular}{c}
  Potpis magistranda\\[2.5em]
  \rule{5cm}{0.4pt}
\end{tabular}
\endgroup

% Zahvalnica — nije obavezna; obrisati % Zahvalnica — nije obavezna; obrisati % Zahvalnica — nije obavezna; obrisati \input{uvodni/zahvalnica} u main.tex ako se ne piše.
\chapter*{ZAHVALNICA}
\thispagestyle{empty}

\napomena{Kratka zahvalnica porodici, mentoru, osoblju fakulteta i komisiji (do pola strane).}

\vspace{1em}
\begin{flushright}
  \RadMjesto, \RadGodina.\\
  \RadAutor
\end{flushright}
 u main.tex ako se ne piše.
\chapter*{ZAHVALNICA}
\thispagestyle{empty}

\napomena{Kratka zahvalnica porodici, mentoru, osoblju fakulteta i komisiji (do pola strane).}

\vspace{1em}
\begin{flushright}
  \RadMjesto, \RadGodina.\\
  \RadAutor
\end{flushright}
 u main.tex ako se ne piše.
\chapter*{ZAHVALNICA}
\thispagestyle{empty}

\napomena{Kratka zahvalnica porodici, mentoru, osoblju fakulteta i komisiji (do pola strane).}

\vspace{1em}
\begin{flushright}
  \RadMjesto, \RadGodina.\\
  \RadAutor
\end{flushright}

% Apstrakt sa ključnim riječima (Smjernice, tačka 4): najviše jedna strana, 5 ključnih riječi.
% Pisati na kraju, kad je rad gotov.
\chapter*{APSTRAKT}
\thispagestyle{empty}

\napomena{Cjeloviti prikaz rada na najviše jednoj strani: predmet istraživanja, metode,
centralna ideja (dvoprolazni asembler za RV32I), rezultati, odgovori na istraživačka
pitanja i naučni doprinos. Pisati na kraju.}

\bigskip
\noindent\underline{\textbf{Ključne riječi}}: asembler, asemblerski prevodilac, RISC-V, RV32I, mašinski kod

% Abstract — istovjetan Apstraktu, na engleskom jeziku.
\begin{otherlanguage}{english}
\chapter*{ABSTRACT}
\thispagestyle{empty}

\napomena{English translation of the Apstrakt (same content, at most one page).}

\bigskip
\noindent\underline{\textbf{Keywords}}: assembler, two-pass assembler, RISC-V, RV32I, machine code
\end{otherlanguage}

\tableofcontents

% --- Centralni dio: brojevi od 1, počevši od UVODA -------------------------
\clearpage
\pagenumbering{arabic}

\nenumerisano{UVOD}
\label{ch:uvod}

\napomena{Kao u primjeru: UVOD je jedinstven tekst bez potpoglavlja, 4--5 strana.
Smjernice: definiše predmet i osnovne pojmove, daje teorijski okvir, ukazuje na nedostatke
postojećih rješenja, navodi istraživačka pitanja i predstavlja strukturu rada.
Teme redom (svaka po jedan ili više pasusa) --- tekst se može preuzeti iz prijave:}

\napomena{1. Kontekst i aktuelnost teme: mjesto asemblera u lancu kompajler → asembler →
linker → loader; širenje RISC-V arhitekture. \parencite{levine2000,hennessy2019}}

\napomena{2. Predmet istraživanja: dvoprolazni asembler za podskup RV32I.}

\napomena{3. Motivacija: razumijevanje veze hardvera i softvera, složenost industrijskih
alata (GNU as, LLVM), obrazovna vrijednost.}

\napomena{4. Cilj i svrha rada: opšti i posebni ciljevi (iz prijave).}

\napomena{5. Istraživačka pitanja: kao u primjeru --- „Prvo istraživačko pitanje glasi: …
Obrazloženje istraživačkog pitanja: …” za glavno i pomoćna pitanja.}

\napomena{6. Metode istraživanja: analiza i sinteza literature, modelovanje/specifikacija,
konstruktivna metoda, eksperimentalna (testiranje), komparativna, studija slučaja.}

\napomena{7. Struktura rada: po jedan pasus za poglavlja
\ref{ch:osnove}--\ref{ch:evaluacija}, zaključak i priloge.}

\chapter{OSNOVE RAČUNARSKIH ARHITEKTURA I ASEMBLERA}
\label{ch:osnove}

\napomena{Teorijsko poglavlje, 8--10 strana. Glavni izvori: \textcite{patterson2020},
\textcite{hennessy2019}, \textcite{harris2022}, \textcite{waterman2019}.}

\section{Arhitektura računarskih sistema}
\label{sec:arhitektura-sistema}
\napomena{Fon Nojmanov model, procesor--memorija--U/I, ciklus dohvati--dekodiraj--izvrši.}

\section{Mašinski jezik i asemblerski jezik}
\label{sec:masinski-asemblerski}
\napomena{Binarni zapis instrukcija; mnemonici, simboličke adrese; odnos 1:1 između
asemblerske i mašinske instrukcije (osim pseudo-instrukcija).}

\section{Instrukcijski skup procesora}
\label{sec:isa}
\napomena{Pojam ISA kao ugovora između hardvera i softvera; vrste instrukcija
(aritmetičke, logičke, memorijske, kontrola toka).}

\section{RISC i CISC arhitekture}
\label{sec:risc-cisc}
\napomena{Istorijat, karakteristike, primjeri (x86 vs.\ ARM/MIPS/RISC-V), load/store princip.}

\section{Pregled RISC-V arhitekture}
\label{sec:risc-v}
\napomena{Nastanak na Berkliju \parencite{asanovic2014,waterman2016}, otvorenost, modularnost
(osnovni skup RV32I/RV64I + ekstenzije M, A, F, D, C), zašto je izabrana za ovaj rad.}

\section{Uloga asemblera u procesu prevođenja programa}
\label{sec:uloga-asemblera}
\napomena{Lanac izvorni kod → kompajler → asembler → objektni fajl → linker → izvršni fajl.
Može dijagram (slika).}

\chapter{TEORIJSKE OSNOVE PREVODILACA}
\label{ch:teorija}

\napomena{Teorijsko poglavlje, 10--12 strana. Uputstvo: „ovdje će trebati istraživanja”.
Glavni izvori: \textcite{aho2006} (pogl.\ 1--6), \textcite{beck1997} (pogl.\ 1--2),
\textcite{cooper2023}, \textcite{nystrom2021}.}

\section{Vrste prevodilaca}
\label{sec:vrste-prevodilaca}
\napomena{Kompajler, interpreter, asembler, (dis)asembler, kros-prevodilac; jednoprolazni
i višeprolazni prevodioci.}

\section{Struktura prevodioca}
\label{sec:struktura-prevodioca}
\napomena{Faze: front-end (analiza) i back-end (sinteza); tabela simbola kao zajednička
struktura. Dijagram faza.}

\section{Leksička analiza}
\label{sec:teorija-leksicka}
\napomena{Tokeni, lekseme, obrasci; regularni izrazi i konačni automati;
ručno pisani skener naspram generatora (lex/flex).}

\section{Sintaksna analiza}
\label{sec:teorija-sintaksna}
\napomena{Kontekstno-slobodne gramatike, BNF/EBNF, stabla parsiranja;
silazno (rekurzivno spuštanje) i uzlazno parsiranje.}

\section{Semantička analiza}
\label{sec:teorija-semanticka}
\napomena{Provjere značenja: definisanost simbola, opsezi vrijednosti, tipovi operanada;
kod asemblera: nedefinisane/duplirane labele, neposredne vrijednosti van opsega.}

\section{Generisanje mašinskog koda}
\label{sec:teorija-generisanje}
\napomena{Kod asemblera: kodiranje instrukcija u bitska polja; dvoprolazni postupak
i reference unaprijed \parencite{beck1997}; jednoprolazni asembler sa „backpatching”-om.}

\section{Obrada grešaka}
\label{sec:teorija-greske}
\napomena{Vrste grešaka (leksičke, sintaksne, semantičke), lokacija greške,
strategije oporavka (panic mode, nastavak na sljedećoj liniji).}

\chapter{SPECIFIKACIJA JEDNOSTAVNE ARHITEKTURE}
\label{ch:specifikacija}

\napomena{Specifikacija ciljne arhitekture: podskup RV32I. 8--10 strana.
Izvor: \textcite{waterman2019}. Mora se poklapati sa \texttt{src/asm/encode.py}.}

\section{Definisanje arhitekture sistema}
\label{sec:definisanje-arhitekture}
\napomena{32-bitna load/store arhitektura; koji dio RV32I se podržava i šta je izostavljeno
(npr.\ \texttt{fence}, \texttt{ecall}, CSR instrukcije).}

\section{Organizacija memorije}
\label{sec:memorija}
\napomena{Bajtno adresiranje, little-endian, poravnanje riječi; raspored sekcija
\texttt{.text} i \texttt{.data} (početne adrese).}

\section{Registri procesora}
\label{sec:registri}
\napomena{32 registra \texttt{x0}--\texttt{x31}, \texttt{x0} uvijek nula, PC;
tabela ABI imena (\texttt{zero}, \texttt{ra}, \texttt{sp}, \texttt{t0}--\texttt{t6}, \texttt{s0}--\texttt{s11}, \texttt{a0}--\texttt{a7}).}

\section{Format instrukcija}
\label{sec:format-instrukcija}
\napomena{Opisati svaki format i raspored bitova neposredne vrijednosti (posebno B i J).}

\begin{table}[H]
  \centering
  \caption{Formati instrukcija RV32I (bitska polja, od bita 31 do bita 0)}
  \label{tab:formati}
  \small
  \begin{tabular}{@{}lllllll@{}}
    \toprule
    Format & 31--25 & 24--20 & 19--15 & 14--12 & 11--7 & 6--0 \\
    \midrule
    R & funct7 & rs2 & rs1 & funct3 & rd & opcode \\
    I & \multicolumn{2}{l}{imm[11:0]} & rs1 & funct3 & rd & opcode \\
    S & imm[11:5] & rs2 & rs1 & funct3 & imm[4:0] & opcode \\
    B & imm[12|10:5] & rs2 & rs1 & funct3 & imm[4:1|11] & opcode \\
    U & \multicolumn{4}{l}{imm[31:12]} & rd & opcode \\
    J & \multicolumn{4}{l}{imm[20|10:1|11|19:12]} & rd & opcode \\
    \bottomrule
  \end{tabular}
\end{table}

\section{Skup podržanih instrukcija}
\label{sec:skup-instrukcija}
\napomena{Kratak pregled po grupama (aritmetičke, logičke, pomjeranja, memorijske, grananja,
skokovi); kompletna tabela opkodova je u Prilogu~\ref{sec:prilog-tabela}.}

\section{Načini adresiranja}
\label{sec:adresiranje}
\napomena{Registarsko, neposredno, bazno sa pomjerajem (\texttt{lw t1, 8(sp)}),
PC-relativno (grane, \texttt{jal}, \texttt{auipc}).}

\section{Primjer izvršavanja instrukcija}
\label{sec:primjer-izvrsavanja}
\napomena{Ručno kodiranje 2--3 instrukcije korak po korak, npr.\
\texttt{add x3, x1, x2} → \texttt{0x002081B3}; jedna I-tip i jedna B-tip instrukcija.}

\chapter{PROJEKTOVANJE ASEMBLERSKOG JEZIKA}
\label{ch:jezik}

\napomena{Specifikacija ulaznog jezika asemblera. 6--8 strana.
Leksička pravila su već implementirana — vidjeti docstring u \texttt{src/asm/lexer.py}.
Izvor za sintaksu: \textcite{riscvasm2024}, \textcite{fsf2024}.}

\section{Sintaksa asemblerskog jezika}
\label{sec:sintaksa}
\napomena{Leksička pravila (komentari \texttt{\#}, identifikatori, brojevi dec/hex/bin,
stringovi) i gramatika u EBNF obliku. Nacrt gramatike (uskladiti sa parserom):}

\begin{lstlisting}[caption={Nacrt gramatike asemblerskog jezika (EBNF)},label={lst:ebnf}]
program     = { line } EOF ;
line        = [ label ] [ statement ] NEWLINE ;
label       = IDENT ":" ;
statement   = directive | instruction ;
directive   = DIRECTIVE [ operands ] ;
instruction = IDENT [ operands ] ;
operands    = operand { "," operand } ;
operand     = IDENT | INTEGER | STRING | INTEGER "(" IDENT ")" ;
\end{lstlisting}

\section{Direktive asemblera}
\label{sec:direktive}
\napomena{Tabela: \texttt{.text}, \texttt{.data}, \texttt{.word}, \texttt{.byte},
\texttt{.ascii}, \texttt{.asciiz}, \texttt{.space} — sintaksa i dejstvo na brojač lokacija.}

\section{Oznake (labeli) i simboli}
\label{sec:labele}
\napomena{Pravila imenovanja, definicija i upotreba, reference unaprijed,
greške: duplirana i nedefinisana labela.}

\section{Podržani tipovi instrukcija}
\label{sec:tipovi-instrukcija}
\napomena{Sintaksa operanada po formatu (R/I/S/B/U/J) i pseudo-instrukcije
(\texttt{nop}, \texttt{li}, \texttt{la}, \texttt{mv}, \texttt{j}, \texttt{ret}) i u šta se prevode.}

\section{Pravila pisanja programa}
\label{sec:pravila}
\napomena{Jedna naredba po liniji, labela i instrukcija u istoj liniji, velika/mala slova,
obavezne sekcije.}

\section{Primjeri asemblerskih programa}
\label{sec:primjeri-programa}
\napomena{2--3 kratka primjera sa objašnjenjem; ostali su u Prilogu~\ref{sec:prilog-primjeri}.}

\lstinputlisting[style=asm,caption={Primjer programa sa petljom (\texttt{examples/loop.asm})},label={lst:loop}]{../../examples/loop.asm}

\chapter{DIZAJN I IMPLEMENTACIJA RAZVIJENOG ASEMBLERSKOG PREVODIOCA}
\label{ch:implementacija}

\napomena{Centralno poglavlje i najveći doprinos rada, 15--20 strana.
Svaka sekcija odgovara jednom modulu u \texttt{src/asm/}. Prikazivati kratke isječke koda,
ne cijele fajlove (cijeli kod je u Prilogu~\ref{sec:prilog-kod}).}

\section{Arhitektura softverskog rješenja}
\label{sec:arhitektura-rjesenja}
\napomena{Dijagram toka: \texttt{.asm} → Lexer → Parser → Pass~1 (tabela simbola, adrese)
→ Pass~2 (kodiranje) → \texttt{.bin}/\texttt{.hex}. Tabela modula i odgovornosti.}

\section{Izbor programskog jezika i razvojnih alata}
\label{sec:alati}
\napomena{Python 3.12 (čitljivost, standardna biblioteka, bez spoljnih zavisnosti),
pytest, Git/GitHub, VS Code/Cursor, \LaTeX. Zašto bez generatora parsera.}

\section{Implementacija leksičkog analizatora}
\label{sec:impl-lexer}
\napomena{\texttt{lexer.py}: vrste tokena, prolaz znak po znak, lokacije (linija/kolona),
sažimanje praznih linija u jedan NEWLINE, dekodiranje brojeva i stringova, greške.}

\section{Implementacija parsera}
\label{sec:impl-parser}
\napomena{\texttt{parser.py}: rekurzivno spuštanje po gramatici iz Listinga~\ref{lst:ebnf};
čvorovi \texttt{LabelDef}, \texttt{Directive}, \texttt{Instruction}.}

\section{Tabela simbola}
\label{sec:impl-simboli}
\napomena{\texttt{symbols.py}: struktura, definisanje i pretraga, detekcija duplikata.}

\section{Prevođenje instrukcija u mašinski kod}
\label{sec:impl-kodiranje}
\napomena{\texttt{encode.py}: tabela opkodova, funkcije za formate, provjera opsega
neposrednih vrijednosti, PC-relativni pomjeraji za grane i skokove, pseudo-instrukcije.}

\section{Generisanje izlaznog fajla}
\label{sec:impl-izlaz}
\napomena{\texttt{output.py}: binarni (little-endian) i heksadecimalni zapis, opciono listing.}

\section{Obrada sintaksnih i semantičkih grešaka}
\label{sec:impl-greske}
\napomena{\texttt{errors.py}: hijerarhija grešaka, format \texttt{fajl:linija:kolona: vrsta: poruka},
primjeri poruka.}

\section{Testiranje funkcionalnosti asemblera}
\label{sec:impl-testiranje}
\napomena{Jedinični testovi po modulu, integracioni testovi nad primjerima, testovi grešaka;
broj testova i pokrivenost.}

\chapter{PRAKTIČNA REALIZACIJA I PRIMJERI RADA}
\label{ch:realizacija}

\napomena{Demonstracija gotovog asemblera na primjerima, 8--10 strana.}

\section{Opis implementiranog sistema}
\label{sec:opis-sistema}
\napomena{Šta asembler podržava (sažetak), kako se pokreće (CLI), ulazi i izlazi.}

\section{Tok prevođenja programa}
\label{sec:tok-prevodjenja}
\napomena{Jedan program kroz sve faze: tokeni → stavke programa → tabela simbola
nakon prvog prolaza → mašinski kod nakon drugog prolaza.}

\section{Primjer jednostavnog programa}
\label{sec:primjer-programa}
\napomena{Npr.\ \texttt{examples/loop.asm} ili program sa \texttt{.data} sekcijom.}

\section{Generisani mašinski kod}
\label{sec:generisani-kod}
\napomena{Sadržaj \texttt{.hex} fajla, objašnjenje 2--3 riječi bit po bit.}

\section{Analiza rada asemblera}
\label{sec:analiza-rada}
\napomena{Primjeri neispravnih programa i poruka o greškama; reference unaprijed.}

\section{Performanse i ograničenja sistema}
\label{sec:performanse}
\napomena{Vrijeme prevođenja za veće ulaze (npr.\ generisan program sa 10\,000 linija),
poznata ograničenja.}

\chapter{EVALUACIJA I DISKUSIJA}
\label{ch:evaluacija}

\napomena{Kritička ocjena rješenja, 5--7 strana. Ovdje se odgovara na pomoćna istraživačka pitanja.}

\section{Prednosti razvijenog rješenja}
\label{sec:prednosti}
\napomena{Jednostavnost, čitljivost, jasne poruke o greškama, bez zavisnosti, testiranost.}

\section{Nedostaci i ograničenja}
\label{sec:nedostaci}
\napomena{Podskup ISA, nema objektnih fajlova/relokacija/linkera, nema makroa, ograničen
skup pseudo-instrukcija.}

\section{Poređenje sa postojećim asemblerima}
\label{sec:poredjenje}
\napomena{Isti programi kroz GNU as (\texttt{riscv64-unknown-elf-as}) — poređenje mašinskog
koda; poređenje poruka o greškama i funkcionalnosti sa RARS i Ripes \parencite{petersen2021}.
Tabela poređenja.}

\section{Mogućnosti unapređenja}
\label{sec:unapredjenja}
\napomena{Ekstenzije M i C, ELF izlaz, makroi, simulator, grafički prikaz faza prevođenja.}

\nenumerisano{ZAKLJUČAK}
\label{ch:zakljucak}

\napomena{Kao u primjeru: ZAKLJUČAK je jedinstven tekst bez potpoglavlja, 2--3 strane.
Smjernice: sublimacija rada i odgovori na istraživačka pitanja; bez novih tvrdnji
i bez citiranja literature. Teme redom:}

\napomena{1. Ostvareni rezultati --- šta je urađeno; eksplicitan odgovor na glavno i
svako pomoćno istraživačko pitanje.}

\napomena{2. Doprinos rada --- obrazovni i praktični.}

\napomena{3. Pravci budućeg razvoja --- kratko, iz Sekcije~\ref{sec:unapredjenja}.}


% --- LITERATURA ----------------------------------------------------------
% \nocite{*} prikazuje SVE izvore iz literatura.bib i dok nisu citirani.
% Ukloniti pred predaju — u spisku ostaju samo citirani izvori.
\nocite{*}
\nenumerisano{LITERATURA}
\printbibliography[nottype=online,heading=subbibliography,title={knjige i članci}]
\printbibliography[type=online,heading=subbibliography,title={Internet izvori}]

% --- PRILOZI ---------------------------------------------------------------
\nenumerisano{PRILOZI}
\label{ch:prilozi}

\prilog{Izvorni kod asemblera}
\label{sec:prilog-kod}
\napomena{Kod se uključuje direktno iz \texttt{src/asm/} — uvijek je ažuran.
Pri kraju rada odlučiti da li ide cio kod ili samo ključni moduli (obim!).}

\lstinputlisting[style=python,caption={\texttt{tokens.py}}]{../../src/asm/tokens.py}
\lstinputlisting[style=python,caption={\texttt{lexer.py}}]{../../src/asm/lexer.py}
\lstinputlisting[style=python,caption={\texttt{parser.py}}]{../../src/asm/parser.py}
\lstinputlisting[style=python,caption={\texttt{ast\_nodes.py}}]{../../src/asm/ast_nodes.py}
\lstinputlisting[style=python,caption={\texttt{symbols.py}}]{../../src/asm/symbols.py}
\lstinputlisting[style=python,caption={\texttt{directives.py}}]{../../src/asm/directives.py}
\lstinputlisting[style=python,caption={\texttt{encode.py}}]{../../src/asm/encode.py}
\lstinputlisting[style=python,caption={\texttt{assembler.py}}]{../../src/asm/assembler.py}
\lstinputlisting[style=python,caption={\texttt{output.py}}]{../../src/asm/output.py}
\lstinputlisting[style=python,caption={\texttt{errors.py}}]{../../src/asm/errors.py}
\lstinputlisting[style=python,caption={\texttt{cli.py}}]{../../src/asm/cli.py}

\prilog{Primjeri ulaznih programa}
\label{sec:prilog-primjeri}
\napomena{Dodati i odgovarajući \texttt{.hex} izlaz za svaki primjer kad asembler bude gotov.}

\lstinputlisting[style=asm,caption={\texttt{examples/add.asm}}]{../../examples/add.asm}
\lstinputlisting[style=asm,caption={\texttt{examples/loop.asm}}]{../../examples/loop.asm}
\lstinputlisting[style=asm,caption={\texttt{examples/memory.asm}}]{../../examples/memory.asm}

\prilog{Tabela instrukcija i opkodova}
\label{sec:prilog-tabela}
\napomena{Nacrt za planirani podskup (izvor: \textcite{waterman2019}); uskladiti sa
\texttt{encode.py} kad se skup instrukcija konačno utvrdi.}

{\small
\begin{longtable}{@{}lllll@{}}
  \caption{Podržane instrukcije RV32I sa opkodovima}\label{tab:opkodovi}\\
  \toprule
  Instrukcija & Format & opcode & funct3 & funct7 \\
  \midrule
  \endfirsthead
  \toprule
  Instrukcija & Format & opcode & funct3 & funct7 \\
  \midrule
  \endhead
  \texttt{add}   & R & 0110011 & 000 & 0000000 \\
  \texttt{sub}   & R & 0110011 & 000 & 0100000 \\
  \texttt{sll}   & R & 0110011 & 001 & 0000000 \\
  \texttt{slt}   & R & 0110011 & 010 & 0000000 \\
  \texttt{sltu}  & R & 0110011 & 011 & 0000000 \\
  \texttt{xor}   & R & 0110011 & 100 & 0000000 \\
  \texttt{srl}   & R & 0110011 & 101 & 0000000 \\
  \texttt{sra}   & R & 0110011 & 101 & 0100000 \\
  \texttt{or}    & R & 0110011 & 110 & 0000000 \\
  \texttt{and}   & R & 0110011 & 111 & 0000000 \\
  \midrule
  \texttt{addi}  & I & 0010011 & 000 & --- \\
  \texttt{xori}  & I & 0010011 & 100 & --- \\
  \texttt{ori}   & I & 0010011 & 110 & --- \\
  \texttt{andi}  & I & 0010011 & 111 & --- \\
  \texttt{lw}    & I & 0000011 & 010 & --- \\
  \texttt{jalr}  & I & 1100111 & 000 & --- \\
  \midrule
  \texttt{sw}    & S & 0100011 & 010 & --- \\
  \midrule
  \texttt{beq}   & B & 1100011 & 000 & --- \\
  \texttt{bne}   & B & 1100011 & 001 & --- \\
  \texttt{blt}   & B & 1100011 & 100 & --- \\
  \texttt{bge}   & B & 1100011 & 101 & --- \\
  \midrule
  \texttt{lui}   & U & 0110111 & --- & --- \\
  \texttt{auipc} & U & 0010111 & --- & --- \\
  \midrule
  \texttt{jal}   & J & 1101111 & --- & --- \\
  \bottomrule
\end{longtable}
}

\prilog{Korisničko uputstvo}
\label{sec:prilog-uputstvo}
\napomena{Instalacija (Python 3.12, \texttt{pip install -e .}), pokretanje, opcije komandne
linije, format izlaza, primjeri poruka o greškama.}

\begin{lstlisting}[caption={Osnovno korišćenje asemblera}]
python -m asm --help
python -m asm examples/add.asm -o out/
\end{lstlisting}


% --- Izjave na kraju rada (ne ulaze u sadržaj) ------------------------------
% Izjava o istovjetnosti štampane i elektronske verzije — na kraju rada, kao u primjeru.
% Predaje se potpisana: zamijeniti skeniranim PDF-om kad bude potpisana.
\clearpage
\thispagestyle{empty}
\begingroup
\setlength{\parindent}{0pt}
\vspace*{1.5cm}
\begin{center}\bfseries Izjava o istovjetnosti štampane i elektronske verzije master rada\end{center}
\vspace{2em}

Ime i prezime autora: \RadAutor\\
Broj indeksa/upisa: \RadIndeks\\
Studijski program: \RadMasterStudije\\
Naslov rada: \RadNaslov\\
Mentor: \RadMentor

Potpisani/a: \textbf{\RadAutor}

\begin{center}\bfseries Izjavljujem\end{center}

da je štampana verzija mog master rada istovjetna elektronskoj verziji koju sam predao/la
za objavljivanje u Digitalni arhiv Univerziteta Crne Gore.

Istovremeno izjavljujem da dozvoljavam objavljivanje mojih ličnih podataka u vezi sa
dobijanjem akademskog naziva master nauka, kao što su ime i prezime, godina i mjesto rođenja,
naslov master rada i datum odbrane rada.

\vspace{3em}
U Podgorici, [dd.mm.gggg.] godine

\vspace{2em}
\hfill\begin{tabular}{c}
  Potpis magistranda\\[2.5em]
  \rule{5cm}{0.4pt}
\end{tabular}
\endgroup

% Izjava o korišćenju (licenca Kreativne zajednice) — na kraju rada, kao u primjeru.
% Predaje se potpisana: zamijeniti skeniranim PDF-om kad bude potpisana.
\clearpage
\thispagestyle{empty}
\begingroup
\setlength{\parindent}{0pt}
\vspace*{1.5cm}
\begin{center}\bfseries IZJAVA O KORIŠĆENJU\end{center}
\vspace{2em}

Ovlašćujem Univerzitetsku biblioteku da u Digitalnom arhivu Univerziteta Crne Gore pohrani
moj master rad pod nazivom:

\begin{center}\bfseries\enquote{\RadNaslov}\end{center}

koji je moje autorsko djelo.

Master rad sa svim prilozima predao/la sam u elektronskom formatu pogodnom za trajno
arhiviranje.

Moj master rad pohranjen u Digitalni arhiv Univerziteta Crne Gore mogu da koriste svi koji
poštuju odredbe sadržane u odabranom tipu licence Kreativne zajednice (\textit{Creative Commons})
za koju sam se odlučio/la.

\begin{enumerate}[leftmargin=1.5cm, itemsep=0pt]
  \item Autorstvo
  \item Autorstvo -- nekomercijalno
  \item Autorstvo -- nekomercijalno -- bez prerade
  \item Autorstvo -- nekomercijalno -- dijeliti pod istim uslovima
  \item Autorstvo -- bez prerade
  \item Autorstvo -- dijeliti pod istim uslovima
\end{enumerate}

\textit{(Molimo da zaokružite samo jednu od šest ponuđenih licenci.)}

\vspace{3em}
U Podgorici, [dd.mm.gggg.] godine

\vspace{2em}
\hfill\begin{tabular}{c}
  Potpis magistranda\\[2.5em]
  \rule{5cm}{0.4pt}
\end{tabular}
\endgroup


\end{document}
